\begin{frame}{Datasets}
  \begin{columns}[T]
    \column{0.42\textwidth}
    \begin{block}{Primary dataset: ScanNet++}
      High-fidelity indoor scenes with registered DSLR images and commodity
      iPhone captures.
    \end{block}

    \begin{block}{Secondary dataset: DTU}
      Controlled object-centric scenes for supplementary evaluation and
      sanity checks in a cleaner setting.
    \end{block}

    \small If compute is limited, we will use curated subsets rather than the
    full datasets.

    \column{0.58\textwidth}
    \centering
    \includegraphics[width=\linewidth]{scannetpp_overview_crop.png}
  \end{columns}

  \vspace{0.08cm}
  \figsource{Yeshwanth et al., ScanNet++, ICCV 2023, Figure 1}
\end{frame}

\begin{frame}{Why ScanNet++?}
  \begin{columns}[T]
    \column{0.63\textwidth}
    \centering
    \includegraphics[width=\linewidth]{scannetpp_table_crop.png}

    \column{0.37\textwidth}
    \begin{itemize}
      \item 460 indoor scenes with high-fidelity geometry
      \item DSLR images support stronger geometric supervision
      \item Commodity captures make sparse-view testing realistic
      \item Dense semantics and NVS setup make the benchmark richer
    \end{itemize}
  \end{columns}

  \vspace{0.08cm}
  \figsource{Yeshwanth et al., ScanNet++, ICCV 2023, Table 1}
\end{frame}

\begin{frame}{Evaluation Plan}
  \small
  \begin{columns}[T]
    \column{0.52\textwidth}
    \begin{itemize}
      \item Reconstruction F-score against available ground-truth geometry
      \item Accuracy, completeness, and Chamfer distance of point clouds
      \item Sparse 2--5 view behavior across scene splits
      \item Limit-specific gates for occlusion and wrong-depth ambiguity
      \item Runtime per scene for practical inference
    \end{itemize}

    \vspace{0.08cm}
    \begin{block}{Success criterion}
      Run 30 must beat the best non-learned baseline overall and improve both
      occlusion and ambiguity subsets by at least the gate margin.
    \end{block}

    \column{0.48\textwidth}
    \centering
    \includegraphics[height=0.46\textheight,keepaspectratio]{mvdust3r_results_table.png}

    \vspace{0.04cm}
    \figsource{Tang et al., MV-DUSt3R+, quantitative benchmark}
  \end{columns}
\end{frame}

\begin{frame}{Final RGB-D Result}
  \scriptsize
  \begin{block}{Run 30 gate decision}
    \texttt{selected\_method = rgbd\_source\_depth\_selected};
    best baseline = \texttt{all\_candidates};
    gate margin = 0.005;
    \texttt{pass\_all\_limits = 1}.
  \end{block}

  \vspace{0.08cm}
  \centering
  \begin{tabular}{lrrrr}
    \toprule
    \textbf{Split/subset} & \textbf{All cand.} & \textbf{Fixed conf.} &
    \textbf{RGB-D selected} & \textbf{Delta vs all} \\
    \midrule
    Val overall & 0.1194 & 0.1123 & 0.1753 & +0.0559 \\
    Val occlusion & 0.1064 & 0.0917 & 0.2522 & +0.1458 \\
    Val ambiguity & 0.1623 & 0.1416 & 0.3000 & +0.1377 \\
    \midrule
    Test overall & 0.1758 & 0.1662 & 0.2764 & +0.1007 \\
    Test occlusion & 0.2111 & 0.1866 & 0.3146 & +0.1035 \\
    Test ambiguity & 0.1621 & 0.1462 & 0.2948 & +0.1327 \\
    \bottomrule
  \end{tabular}
\end{frame}

\begin{frame}{Why RGB-D Works}
  \footnotesize
  \begin{alertblock}{Final wording}
    RGB-only learned extensions did not pass reconstruction-level gates. After
    switching the inference contract to RGB-D, Run 30 uses input source depth
    maps with known camera poses/intrinsics for source-ray correction and passes
    the overall, occlusion, and ambiguity gates on held-out scenes.
  \end{alertblock}

  \begin{itemize}
    \item Source depth anchors each source ray to metric geometry.
    \item Correction preserves candidate count, unlike filtering that hurts
    recall and completeness.
    \item Occlusion and repeated/wrong-depth cases improve because the method
    fixes geometry rather than discarding sparse-view evidence.
  \end{itemize}
\end{frame}
